%----------------------------------------------------------------------------------------
% VECTOREN TEKENEN: SOM EN VEELVOUD. Alle opgaven zijn eigen werk.
%
% Opbouw: Probleem, Start, leerstof, en op een nieuwe pagina de Opdrachten en Test jezelf (geen Koppel in de 2de graad).
% Labels van dit hoofdstuk beginnen met h1-.
% Pakketten: alle pakketten hebben dezelfde leerstof (LPD 13). Pakket 4 en 5 krijgen ook een eigenschap
% bewezen met vectoren (de middenparallel) en twee bewijsopdrachten. Componenten komen pas in hoofdstuk 2
% (pakket 4-5): hier beschrijven we een pijl op een rooster in woorden (3 hokjes naar rechts en 1 omhoog).
% Afspraak in de figuren: een gegeven vector in de kleur van het hoofdstuk (stijl vector), een tweede
% gegeven vector donker (vectorb), de uitkomst dik (som), een verschoven kopie grijs gestreept (hulpvector).
%----------------------------------------------------------------------------------------

\mychapter{1}{Vectoren tekenen}

\begin{probleem}
\vspace{-1mm}\naast{Jens steekt met zijn kajak een rivier over die 60~m breed is. Hij peddelt recht naar de overkant, met een snelheid van 1,5~m/s ten opzichte van het water. Maar het water stroomt, met 0,8~m/s, en neemt hem mee. Waar komt Jens aan de overkant aan? En hoe snel beweegt hij eigenlijk, gezien vanaf de oever?}{\begin{tikzpicture}[scale=.55,font=\scriptsize]
  \fill[accent!8] (-.6,0) rectangle (5.4,3);
  \draw[grijs,thick] (-.6,0) -- (5.4,0) (-.6,3) -- (5.4,3);
  \foreach \y in {.7,1.5,2.3} \draw[lijngrijs,->] (3.4,\y) -- (4.6,\y);
  \fill[inkt] (0,0) circle (2.5pt);
  \draw[vector] (0,0) -- (0,2.1) node[midway,left,text=accent] {1,5 m/s};
  \draw[vectorb] (0,0) -- (1.12,0);
  \node[below,inkt] at (.6,0) {0,8 m/s};
  \draw[grijs,<->] (5.8,0) -- node[right] {60 m} (5.8,3);
  \node[accent] at (2.2,3.35) {?};
\end{tikzpicture}}
\end{probleem}

\lkr{Leerplan \pakketnaam: LPD 13 (vectoren: voorstellen, gelijke en tegengestelde vectoren, de som, het verschil en het scalair veelvoud tekenen, en toepassen op krachten en snelheden).\alleenvoor{4,5}{ In dit pakket ook een eigenschap van de middenparallel, bewezen met vectoren, en twee bewijsopdrachten. Componenten in een assenstelsel, de norm en het ontbinden staan in hoofdstuk 2 (LPD 14).} Voorkennis: verschuivingen (1ste graad), de eigenschappen van een parallellogram, de stelling van Pythagoras (handboek Driehoeksmeting; alleen gebruikt om een tekening na te rekenen). Laat leerlingen op ruitjespapier tekenen, met een geodriehoek. In dit hoofdstuk beschrijven we een pijl op een rooster in woorden (\textit{3 hokjes naar rechts en 1 omhoog}); de notatie met componenten komt\alleenvoor{1,2,3}{ in dit pakket niet aan bod (hoofdstuk 2 staat alleen in pakket 4 en 5)}\alleenvoor{4,5}{ in hoofdstuk 2}.}

Op een weerkaart wijst een pijl aan waar de wind naartoe waait, en hoe langer de pijl, hoe harder het waait. Een getal alleen zegt dat niet: \textit{30~km/h} vertelt niet uit welke hoek de wind komt. Voor zulke grootheden, met een grootte en een richting, gebruik je in de wiskunde en in de fysica \emph{vectoren}. Een verplaatsing, een snelheid en een kracht zijn voorbeelden. In dit hoofdstuk leer je vectoren tekenen, optellen, aftrekken en vermenigvuldigen met een getal. Alles gebeurt met pijlen op een rooster of op schaal, zonder formules. Op het einde los je er het probleem van de kajak mee op.

\begin{oefeningen}[Start]
\oef{h1-start-1}

\oef{h1-verken}
\end{oefeningen}

%========================================================================================
\section{Wat is een vector?}

\subsection{Een verschuiving als pijl}

In de Verken bewoog de robot telkens een aantal hokjes opzij en een aantal hokjes omhoog. Zo'n beweging heet een \emph{verschuiving}. Bij een verschuiving beweegt elk punt evenveel en in dezelfde richting. Verschuif je een driehoek, dan krijg je een even grote driehoek die op dezelfde manier gedraaid staat. Hoe beschrijf je zo'n verschuiving zo kort mogelijk?

\naast{Kijk naar de figuur. Driehoek $ABC$ is verschoven naar driehoek $A'B'C'$. De pijlen van $A$ naar $A'$, van $B$ naar $B'$ en van $C$ naar $C'$ zijn even lang, evenwijdig, en ze wijzen naar dezelfde kant: telkens 5 hokjes naar rechts en 1 hokje omhoog. Eén pijl volstaat dus om de hele verschuiving vast te leggen. Die pijl noemen we een \emph{vector}.}{\begin{tikzpicture}[scale=.5,font=\scriptsize]
  \ruitjes{0}{0}{9}{5}
  \filldraw[fill=accent!12,draw=accent,thick] (1,2) -- (4,1) -- (3,4) -- cycle;
  \filldraw[fill=accent!30,draw=accent,thick] (6,3) -- (9,2) -- (8,5) -- cycle;
  \draw[hulpvector] (1,2) -- (6,3); \draw[hulpvector] (4,1) -- (9,2); \draw[hulpvector] (3,4) -- (8,5);
  \node[left] at (1,2) {$A$}; \node[below] at (4,1) {$B$}; \node[above] at (3,4) {$C$};
  \node[below] at (6,3) {$A'$}; \node[below] at (9,2) {$B'$}; \node[above] at (8,5) {$C'$};
\end{tikzpicture}}

\begin{theorie}
\begin{minipage}[t]{.64\linewidth}
\vspace{0pt}
Een \textbf{\textcolor{accent}{vector}} is een verschuiving. Je stelt hem voor door een pijl. Elke pijl heeft
\begin{itemize}[nosep,leftmargin=5mm]
\item een \textbf{\textcolor{accent}{richting}}: de stand van de rechte waarop de pijl ligt;
\item een \textbf{\textcolor{accent}{zin}}: de kant waarnaar de punt wijst;
\item een \textbf{\textcolor{accent}{grootte}}: de lengte van de pijl.
\end{itemize}
De vector van het \textbf{\textcolor{accent}{beginpunt}} $A$ naar het \textbf{\textcolor{accent}{eindpunt}} $B$ noteer je $\formule{\vect{AB}}$. Een vector met één letter noteer je met een pijl erboven, zoals $\formule{\vect{v}}$.
\end{minipage}\hfill
\begin{minipage}[t]{.32\linewidth}
\vspace{4mm}
\centering
\begin{tikzpicture}[scale=.8,font=\scriptsize]
  \draw[vector] (0,0) -- (2.8,1.2);
  \fill (0,0) circle (1.5pt) node[below] {$A$};
  \fill (2.8,1.2) circle (1.5pt) node[below right] {$B$};
  \node[above left,accent] at (1.4,.6) {$\vect{AB}$};
  \node[grijs,left] at (0,.3) {begin};
  \node[grijs,above] at (2.8,1.3) {eind};
\end{tikzpicture}
\end{minipage}
\end{theorie}

Let op het verschil tussen richting en zin. Een horizontale pijl heeft een horizontale richting, maar hij kan naar rechts of naar links wijzen: dat zijn twee zinnen. Twee pijlen op evenwijdige rechten hebben dezelfde richting. In het dagelijks leven zeggen we \textit{de richting van het station} als we de zin bedoelen; in de wiskunde houden we ze uit elkaar.

\voorbeeld{}{Verschuif de driehoek $PQR$ over de vector $\vect{v}$.\label{vb:h1-verschuif}}{%
\uitleg{De vector $\vect{v}$ gaat 3 hokjes naar rechts en 2 hokjes omhoog. Elk hoekpunt verschuift precies zo. We tellen de hokjes af vanuit $P$, $Q$ en $R$, en verbinden de beelden.}
\naast{\begin{stappenvrij}{5.4cm}{3.6cm}
P \text{ wordt } P' & 3 rechts, 2 omhoog. \\[1ex]
Q \text{ wordt } Q' & Ook 3 rechts, 2 omhoog. \\[1ex]
R \text{ wordt } R' & En nog eens. \\[1ex]
\text{verbind } P', Q', R' & Het beeld is even groot.
\end{stappenvrij}}{\begin{tikzpicture}[scale=.5,font=\scriptsize]
  \ruitjes{0}{0}{9}{5}
  \draw[vector] (0,3) -- (3,5) node[midway,above left,inner sep=1pt] {$\vect{v}$};
  \filldraw[fill=accent!12,draw=accent,thick] (1,0) -- (4,1) -- (2,2) -- cycle;
  \filldraw[fill=accent!30,draw=accent,thick] (4,2) -- (7,3) -- (5,4) -- cycle;
  \draw[hulpvector] (1,0) -- (4,2); \draw[hulpvector] (4,1) -- (7,3); \draw[hulpvector] (2,2) -- (5,4);
  \node[below] at (1,0) {$P$}; \node[below right] at (4,1) {$Q$}; \node[left] at (2,2) {$R$};
  \node[below right] at (4,2) {$P'$}; \node[right] at (7,3) {$Q'$}; \node[above] at (5,4) {$R'$};
\end{tikzpicture}}
De drie grijze pijlen stellen allemaal dezelfde vector $\vect{v}$ voor.}

\begin{oefeningen}
\oef{h1-pijl}
\end{oefeningen}

\subsection{Gelijke vectoren}

In het voorbeeld tekenden we drie pijlen die allemaal dezelfde verschuiving voorstellen. Het zijn drie pijlen, maar één vector. Een vector heeft dus geen vaste plaats: je mag hem tekenen vanuit elk punt dat je wilt. Wanneer stellen twee pijlen dezelfde vector voor?

\begin{theorie}
Twee vectoren zijn \textbf{\textcolor{accent}{gelijk}} als ze dezelfde richting, dezelfde zin en dezelfde grootte hebben. Je schrijft bijvoorbeeld $\formule{\vect{AA'} = \vect{BB'}}$.
\end{theorie}

Twee gelijke vectoren liggen op evenwijdige rechten (of op dezelfde rechte), wijzen naar dezelfde kant en zijn even lang. Is één van die drie anders, dan zijn de vectoren niet gelijk. Een bekende plaats waar gelijke vectoren opduiken, is een parallellogram.

\voorbeeld{}{$ABCD$ is een parallellogram, met $M$ het snijpunt van de diagonalen. Welke vector met letters uit de figuur is gelijk aan $\vect{AB}$? En aan $\vect{AD}$ en aan $\vect{AM}$? Is $\vect{AB} = \vect{CD}$?\label{vb:h1-gelijk}}{%
\uitleg{In een parallellogram zijn de overstaande zijden evenwijdig en even lang, en de diagonalen snijden elkaar middendoor. We kijken telkens ook naar de zin.}
\naast{\begin{stappenvrij}{5.6cm}{4cm}
\vect{AB} = \vect{DC} & Evenwijdig, even lang, allebei naar rechts. \\[1ex]
\vect{AD} = \vect{BC} & Allebei schuin naar boven. \\[1ex]
\vect{AM} = \vect{MC} & $M$ is het midden van $[AC]$. \\[1ex]
\vect{AB} \neq \vect{CD} & $\vect{CD}$ wijst naar links.
\end{stappenvrij}}{\begin{tikzpicture}[scale=.62,font=\scriptsize]
  \coordinate (A) at (0,0); \coordinate (B) at (4,0); \coordinate (C) at (5.4,2.4); \coordinate (D) at (1.4,2.4); \coordinate (M) at (2.7,1.2);
  \fill[accent!8] (A) -- (B) -- (C) -- (D) -- cycle;
  \draw[lijngrijs] (A) -- (C) (B) -- (D);
  \draw[vector] (A) -- (B); \draw[vector] (D) -- (C);
  \draw[vectorb] (A) -- (D); \draw[vectorb] (B) -- (C);
  \fill (M) circle (1.5pt) node[below=2pt] {$M$};
  \node[below left] at (A) {$A$}; \node[below right] at (B) {$B$}; \node[above right] at (C) {$C$}; \node[above left] at (D) {$D$};
\end{tikzpicture}}
$\vect{AB}$ en $\vect{CD}$ hebben dezelfde richting en grootte, maar niet dezelfde zin: ze zijn niet gelijk.}

Wat in het voorbeeld gebeurde, werkt ook omgekeerd: zijn twee overstaande zijden van een vierhoek gelijke vectoren, dan is die vierhoek een parallellogram.

\begin{theorie}
\begin{minipage}[t]{.64\linewidth}
\vspace{0pt}
\begin{stelling}{Vector en parallellogram}
Liggen de punten $A$, $B$, $C$ en $D$ niet op één rechte, dan geldt

\vspace{1mm}
\qquad $\formule{\vect{AB} = \vect{DC} \iff ABCD \text{ is een parallellogram}}$
\end{stelling}
\end{minipage}\hfill
\begin{minipage}[t]{.32\linewidth}
\vspace{2mm}
\centering
\begin{tikzpicture}[scale=.55,font=\scriptsize]
  \coordinate (A) at (0,0); \coordinate (B) at (3.4,0); \coordinate (C) at (4.4,2); \coordinate (D) at (1,2);
  \fill[accent!12] (A) -- (B) -- (C) -- (D) -- cycle;
  \draw[lijngrijs] (A) -- (D) (B) -- (C);
  \draw[vector] (A) -- (B); \draw[vector] (D) -- (C);
  \node[below left] at (A) {$A$}; \node[below right] at (B) {$B$}; \node[above right] at (C) {$C$}; \node[above left] at (D) {$D$};
\end{tikzpicture}
\end{minipage}
\end{theorie}

Waarom klopt dat? Is $\vect{AB} = \vect{DC}$, dan zijn de zijden $[AB]$ en $[DC]$ evenwijdig en even lang. Een vierhoek met twee overstaande zijden die evenwijdig en even lang zijn, is een parallellogram: dat ken je uit de 1ste graad. Let op de volgorde van de letters: $\vect{AB} = \vect{DC}$, en niet $\vect{AB} = \vect{CD}$, want je loopt rond in de volgorde $A$, $B$, $C$, $D$.

\nujij{$PQRS$ is een parallellogram. Welke vector is gelijk aan $\vect{QR}$? En aan $\vect{SR}$?}{$\vect{QR} = \vect{PS}$ en $\vect{SR} = \vect{PQ}$.}

\begin{oefeningen}
\oef{h1-gelijk}
\end{oefeningen}

\subsection{De tegengestelde vector en de nulvector}

Verschuif je een figuur over een vector, hoe krijg je ze dan weer op haar oude plaats? Je verschuift ze even ver terug, in dezelfde richting maar naar de andere kant. Die verschuiving terug heeft een eigen naam.

\begin{theorie}
\begin{minipage}[t]{.64\linewidth}
\vspace{0pt}
De \textbf{\textcolor{accent}{tegengestelde vector}} van $\vect{v}$ is de vector met dezelfde richting en dezelfde grootte, maar met de andere zin. Je noteert hem $\formule{-\vect{v}}$. Zo is $\formule{\vect{BA} = -\vect{AB}}$.
\end{minipage}\hfill
\begin{minipage}[t]{.32\linewidth}
\vspace{1mm}
\centering
\begin{tikzpicture}[scale=.7,font=\scriptsize]
  \draw[vector] (0,0) -- (2.4,1) node[midway,above left,inner sep=1pt] {$\vect{v}$};
  \draw[vectorb] (3.4,1) -- (1,0) node[pos=.4,below right,inner sep=1pt] {$-\vect{v}$};
\end{tikzpicture}
\end{minipage}
\end{theorie}

Ga je eerst van $A$ naar $B$ en dan terug naar $A$, dan ben je niet verschoven. Ook dat is een verschuiving, al is het een vreemde: elk punt blijft op zijn plaats.

\begin{theorie}
De \textbf{\textcolor{accent}{nulvector}} $\formule{\nulvec}$ is de verschuiving waarbij elk punt op zijn plaats blijft. Zo is $\formule{\vect{AA} = \nulvec}$.
\end{theorie}

De nulvector is een pijl met lengte nul: een punt. Hij heeft grootte 0, en geen richting of zin.

\voorbeeld{}{In de figuur is $ABCDEF$ een regelmatige zeshoek met middelpunt $M$. Zoek in de figuur twee pijlen die de tegengestelde vector van $\vect{MA}$ voorstellen. Welke vector is $\vect{MA} + \vect{AM}$?\label{vb:h1-tegen}}{%
\uitleg{De tegengestelde van $\vect{MA}$ is even lang, evenwijdig en wijst de andere kant op. Een regelmatige zeshoek bestaat uit zes gelijkzijdige driehoeken, dus alle stukken in de figuur zijn even lang.}
\naast{\begin{stappenvrij}{5.4cm}{4.4cm}
-\vect{MA} = \vect{AM} & Dezelfde pijl, omgekeerd. \\[1ex]
-\vect{MA} = \vect{MD} & $D$ ligt recht tegenover $A$. \\[1ex]
-\vect{MA} = \vect{BC} = \vect{FE} & Ook evenwijdig met $[AD]$. \\[1ex]
\vect{MA} + \vect{AM} = \nulvec & Heen en terug.
\end{stappenvrij}}{\begin{tikzpicture}[scale=.75,font=\scriptsize]
  \foreach \i/\n in {0/D,60/C,120/B,180/A,240/F,300/E} \coordinate (\n) at (\i:1.6);
  \coordinate (M) at (0,0);
  \fill[accent!8] (A) -- (B) -- (C) -- (D) -- (E) -- (F) -- cycle;
  \draw[lijngrijs] (A) -- (B) -- (C) -- (D) -- (E) -- (F) -- cycle (B) -- (E) (C) -- (F);
  \draw[vector] (M) -- (A);
  \draw[vectorb] (M) -- (D);
  \fill (M) circle (1.5pt) node[below=2pt] {$M$};
  \foreach \n/\p in {A/left,B/above left,C/above right,D/right,E/below right,F/below left} \node[\p] at (\n) {$\n$};
\end{tikzpicture}}
Er zijn dus meer dan twee goede antwoorden: elke pijl van lengte $\lengte{MA}$ die horizontaal naar rechts wijst.}

\begin{oefeningen}
\oef{h1-tegen}
\end{oefeningen}

%========================================================================================
\section{De som van vectoren}

\subsection{Twee verschuivingen na elkaar}

In de Verken voerde de robot twee opdrachten na elkaar uit, en kwam uit op een plaats die hij ook met één beweging kon bereiken. Twee verschuivingen na elkaar geven weer een verschuiving. Die noemen we de som van de twee vectoren. Hoe teken je die som?

\begin{theorie}
\begin{minipage}[t]{.6\linewidth}
\vspace{0pt}
De \textbf{\textcolor{accent}{som}} $\formule{\vect{u} + \vect{v}}$ is de verschuiving over $\vect{u}$ gevolgd door de verschuiving over $\vect{v}$.

Je tekent ze met de \textbf{\textcolor{accent}{kop-staartregel}}: teken $\vect{v}$ met zijn staart (beginpunt) aan de kop (de punt) van $\vect{u}$. De som gaat van de staart van $\vect{u}$ naar de kop van $\vect{v}$.
\end{minipage}\hfill
\begin{minipage}[t]{.36\linewidth}
\vspace{2mm}
\centering
\begin{tikzpicture}[scale=.7,font=\scriptsize]
  \draw[vector] (0,0) -- (2.6,.6) node[midway,below right,inner sep=1pt] {$\vect{u}$};
  \draw[vectorb] (2.6,.6) -- (3.4,2.4) node[midway,right] {$\vect{v}$};
  \draw[som] (0,0) -- (3.4,2.4) node[midway,above left,inner sep=1pt] {$\vect{u} + \vect{v}$};
\end{tikzpicture}
\end{minipage}
\end{theorie}

Het idee: wie eerst over $\vect{u}$ en dan over $\vect{v}$ verschuift, komt op dezelfde plaats uit als wie in één keer over $\vect{u} + \vect{v}$ verschuift. Daarom zit de staart van $\vect{v}$ aan de kop van $\vect{u}$: de tweede verschuiving begint waar de eerste eindigt.

\voorbeeld{}{Teken de som van de vectoren $\vect{u}$ en $\vect{v}$ uit de linkse figuur.\label{vb:h1-kopstaart}}{%
\uitleg{$\vect{u}$ gaat 4 hokjes naar rechts en 2 omhoog, $\vect{v}$ gaat 2 naar rechts en 3 omlaag. We tekenen $\vect{u}$ vanuit een punt $P$, en $\vect{v}$ vanaf de kop van $\vect{u}$.}
\naast{\begin{stappenvrij}{4.4cm}{3.8cm}
\text{teken } \vect{u} \text{ vanuit } P & 4 rechts, 2 omhoog. \\[1ex]
\text{teken } \vect{v} \text{ aan de kop} & 2 rechts, 3 omlaag. \\[1ex]
\text{som: van } P \text{ naar de kop} & 6 rechts, 1 omlaag.
\end{stappenvrij}}{\begin{tikzpicture}[scale=.5,font=\scriptsize]
  \ruitjes{0}{0}{4}{5}
  \draw[vector] (0,0) -- (4,2) node[midway,below right,inner sep=1pt] {$\vect{u}$};
  \draw[vectorb] (1,5) -- (3,2) node[midway,right] {$\vect{v}$};
  \begin{scope}[xshift=5.5cm]
  \ruitjes{0}{0}{7}{5}
  \draw[vector] (0,2) -- (4,4) node[midway,above left,inner sep=1pt] {$\vect{u}$};
  \draw[vectorb] (4,4) -- (6,1) node[midway,right] {$\vect{v}$};
  \draw[som] (0,2) -- (6,1);
  \node[below,accent] at (3,1.5) {$\vect{u} + \vect{v}$};
  \fill (0,2) circle (2.5pt) node[left] {$P$};
  \end{scope}
\end{tikzpicture}}
De som gaat 6 hokjes naar rechts en 1 hokje omlaag. Merk op: de hokjes naar rechts tel je op ($4 + 2$), en 2 omhoog en 3 omlaag geeft samen 1 omlaag.}

\begin{werkwijze}
Teken de eerste vector vanuit een punt naar keuze. Teken de tweede vector met zijn staart aan de kop van de eerste: neem hem over met dezelfde richting, zin en grootte. Teken de som van de staart van de eerste naar de kop van de tweede.
\end{werkwijze}

Een veelgemaakte fout is de som tekenen van kop naar staart, dus in de verkeerde zin. Controleer altijd: de som begint waar de eerste vector begint.

\begin{oefeningen}
\oef{h1-kopstaart}
\end{oefeningen}

Liggen de pijlen al kop aan staart tussen punten met een naam, dan heb je geen tekening meer nodig. Wie van $A$ naar $B$ gaat en dan van $B$ naar $C$, komt in $C$ uit.

\begin{theorie}
\begin{stelling}{Betrekking van Chasles}
Voor elke drie punten $A$, $B$ en $C$ geldt

\vspace{1mm}
\qquad $\formule{\vect{AB} + \vect{BC} = \vect{AC}}$
\end{stelling}
\end{theorie}

De regel is genoemd naar de Franse wiskundige Michel Chasles (1793-1880). Hij werkt alleen als het eindpunt van de eerste vector het beginpunt van de tweede is: de middelste letters moeten gelijk zijn. Je mag hem ook meer dan één keer na elkaar gebruiken.

\voorbeeld{}{Schrijf als één vector: $\vect{PQ} + \vect{QR} + \vect{RS}$.\label{vb:h1-chasles}}{%
\uitleg{De letters sluiten telkens op elkaar aan: $Q$ en $Q$, $R$ en $R$. We gebruiken de betrekking van Chasles twee keer.}
\begin{stappen}
\vect{PQ} + \vect{QR} + \vect{RS} = \vect{PR} + \vect{RS} & Chasles voor de eerste twee. \\[1ex]
\phantom{\vect{PQ} + \vect{QR} + \vect{RS}} = \vect{PS} & En nog eens.
\end{stappen}
Van $P$ naar $Q$, naar $R$, naar $S$: je komt van $P$ in $S$ uit.}

\voorbeeld{}{Schrijf als één vector: $\vect{BC} + \vect{AB}$.\label{vb:h1-chasles2}}{%
\uitleg{De letters sluiten niet aan: $C$ en $A$ zijn verschillend. Maar in de andere volgorde sluiten ze wel aan, en de volgorde van een som mag je wisselen (dat zie je straks bij de parallellogramregel).}
\begin{stappen}
\vect{BC} + \vect{AB} = \vect{AB} + \vect{BC} & Wissel de volgorde. \\[1ex]
\phantom{\vect{BC} + \vect{AB}} = \vect{AC} & Chasles.
\end{stappen}}

\begin{oefeningen}
\oef{h1-chasles}
\end{oefeningen}

\subsection{De parallellogramregel}

Twee sleepboten trekken een schip de haven uit. Beide kabels zitten vast aan de neus van het schip: de twee krachten vertrekken uit hetzelfde punt. De kop-staartregel past dan niet meteen, want de pijlen liggen niet kop aan staart. Hoe vind je toch de som?

\naast{Het idee: verschuif de tweede pijl, zodat hij aan de kop van de eerste begint. Dat mag, want een vector heeft geen vaste plaats. Doe je hetzelfde met de eerste pijl aan de kop van de tweede, dan ontstaat een parallellogram. De som is de diagonaal van dat parallellogram die in het gemeenschappelijke beginpunt vertrekt.}{\begin{tikzpicture}[scale=.62,font=\scriptsize]
  \draw[vector] (0,0) -- (3,.5) node[midway,below,inner sep=2pt] {$\vect{u}$};
  \draw[vectorb] (0,0) -- (1,2) node[midway,left] {$\vect{v}$};
  \draw[hulpvector] (3,.5) -- (4,2.5);
  \draw[hulpvector] (1,2) -- (4,2.5);
  \draw[som] (0,0) -- (4,2.5) node[pos=.62,above left,inner sep=1pt] {$\vect{u} + \vect{v}$};
  \fill (0,0) circle (1.8pt) node[below left] {$O$};
\end{tikzpicture}}

\begin{theorie}
\begin{minipage}[t]{.64\linewidth}
\vspace{0pt}
\begin{stelling}{Parallellogramregel}
Teken $\vect{u}$ en $\vect{v}$ vanuit hetzelfde punt $O$, en vul ze aan tot een parallellogram. De som $\formule{\vect{u} + \vect{v}}$ is de diagonaal van dat parallellogram vanuit $O$.
\end{stelling}
\end{minipage}\hfill
\begin{minipage}[t]{.32\linewidth}
\vspace{1mm}
\centering
\begin{tikzpicture}[scale=.55,font=\scriptsize]
  \fill[accent!10] (0,0) -- (3,.5) -- (4,2.5) -- (1,2) -- cycle;
  \draw[lijngrijs] (3,.5) -- (4,2.5) -- (1,2);
  \draw[vector] (0,0) -- (3,.5) node[midway,below,inner sep=2pt] {$\vect{u}$};
  \draw[vectorb] (0,0) -- (1,2) node[midway,left] {$\vect{v}$};
  \draw[som] (0,0) -- (4,2.5);
  \node[below left] at (0,0) {$O$};
\end{tikzpicture}
\end{minipage}
\end{theorie}

Waarom klopt dat? De zijde tegenover $\vect{v}$ in het parallellogram stelt dezelfde vector $\vect{v}$ voor: ze is evenwijdig, even lang en wijst naar dezelfde kant. Langs die zijde liggen $\vect{u}$ en $\vect{v}$ kop aan staart, dus de kop-staartregel geeft dezelfde diagonaal. De twee regels geven altijd dezelfde som: kies de regel die het best past bij de figuur.

\voorbeeld{}{Twee touwen trekken aan een paal in het punt $O$. De krachten zijn de vectoren $\vect{a}$ en $\vect{b}$ in de figuur. Teken de som met de parallellogramregel.\label{vb:h1-parallel}}{%
\uitleg{Beide vectoren vertrekken in $O$. We tekenen door de kop van $\vect{a}$ een pijl gelijk aan $\vect{b}$, en door de kop van $\vect{b}$ een pijl gelijk aan $\vect{a}$. Waar ze samenkomen, ligt de vierde hoek.}
\naast{\begin{stappenvrij}{5.6cm}{2.8cm}
\vect{a}: 4 \text{ rechts}, 1 \text{ omhoog} & Kop van $\vect{a}$. \\[1ex]
\vect{b}: 1 \text{ links}, 3 \text{ omhoog} & Kop van $\vect{b}$. \\[1ex]
\text{vierde hoek: } 3 \text{ rechts}, 4 \text{ omhoog} & Vanuit $O$.
\end{stappenvrij}}{\begin{tikzpicture}[scale=.5,font=\scriptsize]
  \ruitjes{-2}{0}{5}{5}
  \fill[accent!10] (0,0) -- (4,1) -- (3,4) -- (-1,3) -- cycle;
  \draw[hulpvector] (4,1) -- (3,4); \draw[hulpvector] (-1,3) -- (3,4);
  \draw[vector] (0,0) -- (4,1) node[midway,below right,inner sep=1pt] {$\vect{a}$};
  \draw[vectorb] (0,0) -- (-1,3) node[midway,left] {$\vect{b}$};
  \draw[som] (0,0) -- (3,4);
  \node[accent,right] at (1.9,2.2) {$\vect{a} + \vect{b}$};
  \fill (0,0) circle (2.5pt) node[below] {$O$};
\end{tikzpicture}}
De som gaat 3 hokjes naar rechts en 4 omhoog. Op ruitjespapier is de som ook $4 - 1 = 3$ naar rechts en $1 + 3 = 4$ omhoog.}

\begin{werkwijze}
Liggen de vectoren met hun beginpunt in hetzelfde punt, gebruik dan de parallellogramregel: teken door de kop van elke vector een pijl gelijk aan de andere, en teken de diagonaal vanuit het gemeenschappelijke beginpunt. Liggen ze kop aan staart, gebruik dan de kop-staartregel.
\end{werkwijze}

\begin{oefeningen}
\oef{h1-parallel}
\end{oefeningen}

\subsection{Wisselen mag}

Maakt de volgorde uit? In de Verken kwam de robot op dezelfde plaats uit, of hij nu eerst opdracht A of eerst opdracht B uitvoerde. Het parallellogram toont waarom dat altijd zo is.

\begin{theorie}
\begin{minipage}[t]{.64\linewidth}
\vspace{0pt}
\begin{stelling}{De som is commutatief}
Voor alle vectoren $\vect{u}$ en $\vect{v}$ geldt

\vspace{1mm}
\qquad $\formule{\vect{u} + \vect{v} = \vect{v} + \vect{u}}$
\end{stelling}
\end{minipage}\hfill
\begin{minipage}[t]{.32\linewidth}
\vspace{1mm}
\centering
\begin{tikzpicture}[scale=.55,font=\scriptsize]
  \draw[vector] (0,0) -- (3,.5) node[midway,below,inner sep=2pt] {$\vect{u}$};
  \draw[vectorb] (3,.5) -- (4,2.5) node[midway,right] {$\vect{v}$};
  \draw[vectorb] (0,0) -- (1,2) node[midway,left] {$\vect{v}$};
  \draw[vector] (1,2) -- (4,2.5) node[midway,above,inner sep=2pt] {$\vect{u}$};
  \draw[som] (0,0) -- (4,2.5);
\end{tikzpicture}
\end{minipage}
\end{theorie}

In het parallellogram kun je op twee manieren van het beginpunt naar de overstaande hoek: onderaan eerst $\vect{u}$ en dan $\vect{v}$, of bovenaan eerst $\vect{v}$ en dan $\vect{u}$. Je komt in dezelfde hoek uit. Dat is ook wat je verwacht: tien stappen naar rechts en dan vijf naar voor brengt je op dezelfde plaats als eerst vijf naar voor en dan tien naar rechts.

Bij drie of meer vectoren ga je verder met de kop-staartregel: elke volgende vector begint aan de kop van de vorige. De som gaat van de staart van de eerste naar de kop van de laatste. Ook hier mag je de volgorde kiezen.

\voorbeeld{}{Teken de som $\vect{u} + \vect{v} + \vect{w}$ van de drie vectoren in de figuur.\label{vb:h1-drie}}{%
\uitleg{We tekenen de drie vectoren kop aan staart, vanuit een punt $P$: eerst $\vect{u}$ (3 rechts, 1 omhoog), dan $\vect{v}$ (1 rechts, 3 omhoog), dan $\vect{w}$ (3 links, 1 omhoog).}
\naast{\begin{stappenvrij}{4.4cm}{3.4cm}
\text{kop-staart: } \vect{u}, \vect{v}, \vect{w} & Telkens aan de kop. \\[1ex]
\text{som: van } P \text{ naar de kop} & 1 rechts, 5 omhoog.
\end{stappenvrij}}{\begin{tikzpicture}[scale=.5,font=\scriptsize]
  \ruitjes{0}{0}{5}{5}
  \draw[vector] (0,0) -- (3,1) node[midway,above left,inner sep=1pt] {$\vect{u}$};
  \draw[vectorb] (4,1) -- (5,4) node[midway,left] {$\vect{v}$};
  \draw[vector,dashed] (3,4) -- (0,5) node[midway,below,inner sep=2pt] {$\vect{w}$};
  \begin{scope}[xshift=6.5cm]
  \ruitjes{0}{0}{4}{5}
  \draw[vector] (0,0) -- (3,1); \draw[vectorb] (3,1) -- (4,4); \draw[vector,dashed] (4,4) -- (1,5);
  \draw[som] (0,0) -- (1,5);
  \node[accent,left] at (.5,2.6) {som};
  \fill (0,0) circle (2.5pt) node[left] {$P$};
  \end{scope}
\end{tikzpicture}}
De som gaat 1 hokje naar rechts en 5 omhoog. Probeer zelf een andere volgorde, bijvoorbeeld $\vect{w} + \vect{u} + \vect{v}$: je komt in hetzelfde punt uit.}

Komt de kop van de laatste vector terug in het beginpunt, dan is de som de nulvector: de pijlen vormen dan een gesloten veelhoek. Dat zie je ook in het voorbeeld van de zeshoek: $\vect{AB} + \vect{BC} + \vect{CD} + \vect{DE} + \vect{EF} + \vect{FA} = \vect{AA} = \nulvec$.

\begin{oefeningen}
\oef{h1-drie}
\end{oefeningen}

\subsection{Het verschil van twee vectoren}

Bij getallen is aftrekken hetzelfde als het tegengestelde optellen: $7 - 3 = 7 + ({-3})$. Met vectoren spreken we hetzelfde af. Zo heb je geen nieuwe tekenregel nodig.

\begin{theorie}
Het \textbf{\textcolor{accent}{verschil}} van $\vect{u}$ en $\vect{v}$ is de som van $\vect{u}$ en de tegengestelde van $\vect{v}$:

\vspace{1mm}
\qquad $\formule{\vect{u} - \vect{v} = \vect{u} + (-\vect{v})}$
\end{theorie}

Er is ook een andere manier om naar het verschil te kijken. Het verschil $\vect{u} - \vect{v}$ is de vector die je bij $\vect{v}$ moet optellen om $\vect{u}$ te krijgen, zoals $7 - 3$ het getal is dat je bij 3 optelt om 7 te krijgen.

\voorbeeld{}{Teken $\vect{u} - \vect{v}$ voor de vectoren in de figuur.\label{vb:h1-verschil}}{%
\uitleg{We tellen bij $\vect{u}$ de tegengestelde van $\vect{v}$ op. $\vect{u}$ gaat 4 rechts en 1 omhoog; $\vect{v}$ gaat 1 rechts en 2 omhoog, dus $-\vect{v}$ gaat 1 links en 2 omlaag.}
\naast{\begin{stappenvrij}{4.4cm}{3.4cm}
\text{teken } \vect{u} \text{ vanuit } P & 4 rechts, 1 omhoog. \\[1ex]
\text{teken } -\vect{v} \text{ aan de kop} & 1 links, 2 omlaag. \\[1ex]
\text{verschil: van } P \text{ naar de kop} & 3 rechts, 1 omlaag.
\end{stappenvrij}}{\begin{tikzpicture}[scale=.5,font=\scriptsize]
  \ruitjes{0}{0}{4}{4}
  \draw[vector] (0,2) -- (4,3) node[midway,above left,inner sep=1pt] {$\vect{u}$};
  \draw[vectorb] (1,0) -- (2,2) node[midway,right] {$\vect{v}$};
  \begin{scope}[xshift=5.5cm]
  \ruitjes{0}{0}{4}{4}
  \draw[vector] (0,2) -- (4,3) node[midway,above left,inner sep=1pt] {$\vect{u}$};
  \draw[vectorb] (4,3) -- (3,1) node[midway,right] {$-\vect{v}$};
  \draw[som] (0,2) -- (3,1) node[midway,below left,inner sep=1pt] {$\vect{u} - \vect{v}$};
  \fill (0,2) circle (2.5pt) node[left] {$P$};
  \end{scope}
\end{tikzpicture}}
Controle: tel je $\vect{v}$ (1 rechts, 2 omhoog) op bij $\vect{u} - \vect{v}$ (3 rechts, 1 omlaag), dan krijg je 4 rechts en 1 omhoog, dus $\vect{u}$.}

Tekenen de vectoren vanuit hetzelfde punt, dan is er een snellere manier. Neem $\vect{OA}$ en $\vect{OB}$. Welke vector moet je bij $\vect{OA}$ optellen om $\vect{OB}$ te krijgen? Volgens Chasles is dat $\vect{AB}$, want $\vect{OA} + \vect{AB} = \vect{OB}$.

\begin{theorie}
\begin{minipage}[t]{.64\linewidth}
\vspace{0pt}
\begin{stelling}{Verschil vanuit één punt}
Voor elke drie punten $O$, $A$ en $B$ geldt

\vspace{1mm}
\qquad $\formule{\vect{OB} - \vect{OA} = \vect{AB}}$
\end{stelling}
\end{minipage}\hfill
\begin{minipage}[t]{.32\linewidth}
\vspace{1mm}
\centering
\begin{tikzpicture}[scale=.55,font=\scriptsize]
  \draw[vector] (0,0) -- (3.4,.6) node[midway,below,inner sep=2pt] {$\vect{OA}$};
  \draw[vectorb] (0,0) -- (1.6,2.4) node[midway,left] {$\vect{OB}$};
  \draw[som] (3.4,.6) -- (1.6,2.4);
  \node[accent,above right,inner sep=1pt] at (2.5,1.5) {$\vect{AB}$};
  \node[below left] at (0,0) {$O$}; \node[right] at (3.4,.6) {$A$}; \node[above] at (1.6,2.4) {$B$};
\end{tikzpicture}
\end{minipage}
\end{theorie}

Let op de volgorde: $\vect{OB} - \vect{OA}$ is de pijl van $A$ naar $B$, dus van de kop van de vector die je aftrekt naar de kop van de andere. Deze eigenschap komt nog vaak terug: ze zegt dat je elke vector kunt schrijven als een verschil van twee vectoren vanuit een vast punt.

\voorbeeld{}{$ABCD$ is een parallellogram. Schrijf $\vect{AB} - \vect{AD}$ als één vector met letters uit de figuur.\label{vb:h1-verschil2}}{%
\uitleg{Beide vectoren vertrekken in $A$. Het verschil is de pijl van de kop van $\vect{AD}$ naar de kop van $\vect{AB}$.}
\naast{\begin{stappenvrij}{5cm}{4cm}
\vect{AB} - \vect{AD} = \vect{DB} & Verschil vanuit één punt.
\end{stappenvrij}}{\begin{tikzpicture}[scale=.5,font=\scriptsize]
  \coordinate (A) at (0,0); \coordinate (B) at (4,0); \coordinate (C) at (5.2,2.2); \coordinate (D) at (1.2,2.2);
  \fill[accent!8] (A) -- (B) -- (C) -- (D) -- cycle;
  \draw[lijngrijs] (B) -- (C) -- (D);
  \draw[vector] (A) -- (B); \draw[vectorb] (A) -- (D);
  \draw[som] (D) -- (B);
  \node[below left] at (A) {$A$}; \node[below right] at (B) {$B$}; \node[above right] at (C) {$C$}; \node[above left] at (D) {$D$};
\end{tikzpicture}}
De ene diagonaal van het parallellogram is de som $\vect{AB} + \vect{AD} = \vect{AC}$, de andere is het verschil $\vect{AB} - \vect{AD} = \vect{DB}$.}

\begin{oefeningen}
\oef{h1-verschil}
\end{oefeningen}

%========================================================================================
\section{Een vector vermenigvuldigen met een getal}

\subsection{Het veelvoud van een vector}

Wat is $\vect{v} + \vect{v}$? Met de kop-staartregel liggen de twee pijlen in elkaars verlengde: de som heeft dezelfde richting en zin als $\vect{v}$, en is twee keer zo lang. Net zoals bij getallen schrijven we $2\vect{v}$. Op dezelfde manier is $\vect{v} + \vect{v} + \vect{v} = 3\vect{v}$, drie keer zo lang. En wat zou $\tfrac12\vect{v}$ zijn? Een vector die, twee keer opgeteld, $\vect{v}$ geeft: half zo lang. Zo komen we tot een vermenigvuldiging met elk getal.

\begin{theorie}
\begin{minipage}[t]{.62\linewidth}
\vspace{0pt}
Voor een getal $k$ en een vector $\vect{v}$ is het \textbf{\textcolor{accent}{veelvoud}} $\formule{k\,\vect{v}}$ de vector
\begin{itemize}[nosep,leftmargin=5mm]
\item met dezelfde richting als $\vect{v}$;
\item met een grootte die $\abs{k}$ keer de grootte van $\vect{v}$ is;
\item met dezelfde zin als $\vect{v}$ als $k > 0$, en de andere zin als $k < 0$.
\end{itemize}
Verder is $\formule{0\,\vect{v} = \nulvec}$.
\end{minipage}\hfill
\begin{minipage}[t]{.34\linewidth}
\vspace{1mm}
\centering
\begin{tikzpicture}[scale=.5,font=\scriptsize]
  \draw[vector] (0,2.4) -- (2,3) node[midway,above left,inner sep=1pt] {$\vect{v}$};
  \draw[vector] (0,1.2) -- (4,2.4) node[pos=.7,above left,inner sep=1pt] {$2\vect{v}$};
  \draw[vectorb] (3,1) -- (0,.1) node[midway,below right,inner sep=1pt] {${-1{,}5}\,\vect{v}$};
\end{tikzpicture}
\end{minipage}
\end{theorie}

Het getal $k$ heet in deze context een \textit{scalair}: het schaalt de vector. Een negatief getal draait de zin om. Zo is $({-1})\,\vect{v} = -\vect{v}$, de tegengestelde vector. Let op: $\abs{k}$ is de absolute waarde, dus $-2\vect{v}$ is ook twee keer zo lang als $\vect{v}$, maar wijst de andere kant op.

\voorbeeld{}{Teken $2\vect{v}$, $-\vect{v}$, $\tfrac12\vect{v}$ en $-\tfrac32\vect{v}$ voor de vector $\vect{v}$ in de figuur.\label{vb:h1-veelvoud}}{%
\uitleg{$\vect{v}$ gaat 2 hokjes naar rechts en 4 omhoog. Vermenigvuldigen met $k$ vermenigvuldigt beide aantallen hokjes met $k$; een negatief getal keert de zin om.}
\naast{\begin{stappenvrij}{4.6cm}{3.4cm}
2\vect{v}: 4 \text{ rechts}, 8 \text{ omhoog} & Twee keer zo lang. \\[1ex]
-\vect{v}: 2 \text{ links}, 4 \text{ omlaag} & Andere zin. \\[1ex]
\tfrac12\vect{v}: 1 \text{ rechts}, 2 \text{ omhoog} & Half zo lang. \\[1ex]
-\tfrac32\vect{v}: 3 \text{ links}, 6 \text{ omlaag} & Anderhalf keer, andere zin.
\end{stappenvrij}}{\begin{tikzpicture}[scale=.36,font=\scriptsize]
  \ruitjes{0}{0}{15}{8}
  \draw[vector] (0,2) -- (2,6) node[midway,left] {$\vect{v}$};
  \draw[vector] (3,0) -- (7,8) node[midway,left] {$2\vect{v}$};
  \draw[vectorb] (9,6) -- (7,2) node[midway,right] {$-\vect{v}$};
  \draw[vector] (10,1) -- (11,3) node[midway,left] {$\tfrac12\vect{v}$};
  \draw[vectorb] (15,8) -- (12,2) node[midway,right] {$-\tfrac32\vect{v}$};
\end{tikzpicture}}
Alle pijlen zijn evenwijdig: een veelvoud verandert nooit de richting.}

\begin{oefeningen}
\oef{h1-veelvoud}
\end{oefeningen}

Omdat een veelvoud de richting niet verandert, zijn $\vect{v}$ en $k\,\vect{v}$ altijd evenwijdig. Omgekeerd kun je een vector die evenwijdig is met $\vect{v}$ (en niet de nulvector is) altijd schrijven als een veelvoud van $\vect{v}$. Liggen de pijlen op een rechte, dan lees je het getal $k$ af door de lengtes te vergelijken.

\voorbeeld{}{De punten $A$, $B$, $C$ en $M$ liggen op een rechte, zoals op de figuur. $M$ is het midden van $[AB]$. Vul aan: $\vect{AB} = \ldots\,\vect{AM}$, $\vect{AM} = \ldots\,\vect{AB}$, $\vect{BA} = \ldots\,\vect{AM}$ en $\vect{AC} = \ldots\,\vect{AM}$.\label{vb:h1-factor}}{%
\uitleg{We vergelijken de lengtes: $\lengte{AM} = 2$, $\lengte{AB} = 4$ en $\lengte{AC} = 6$ eenheden. Wijzen twee vectoren naar dezelfde kant, dan is $k$ positief; anders negatief.}
\naast{\begin{stappenvrij}{4.6cm}{4.4cm}
\vect{AB} = 2\,\vect{AM} & Twee keer zo lang, dezelfde zin. \\[1ex]
\vect{AM} = \tfrac12\,\vect{AB} & Half zo lang. \\[1ex]
\vect{BA} = -2\,\vect{AM} & Twee keer zo lang, andere zin. \\[1ex]
\vect{AC} = 3\,\vect{AM} & Drie keer zo lang.
\end{stappenvrij}}{\begin{tikzpicture}[scale=.5,font=\scriptsize]
  \draw[grijs] (-.5,0) -- (6.5,0);
  \foreach \x in {0,1,...,6} \draw[grijs] (\x,.1) -- (\x,-.1);
  \foreach \x/\n in {0/A,2/M,4/B,6/C} {\fill (\x,0) circle (2.5pt); \node[below=2pt] at (\x,0) {$\n$};}
\end{tikzpicture}}
In het bijzonder: $M$ is het midden van $[AB]$ precies dan als $\vect{AM} = \tfrac12\,\vect{AB}$.}

\begin{oefeningen}
\oef{h1-factor}
\end{oefeningen}

\subsection{Som en veelvoud samen}

Met de som en het veelvoud samen kun je nieuwe vectoren maken, zoals $2\vect{u} + \vect{v}$ of $\vect{u} - 3\vect{v}$. Je tekent eerst de veelvouden, en telt ze dan op met de kop-staartregel.

\voorbeeld{}{Teken $2\vect{u} - \vect{v}$ voor de vectoren $\vect{u}$ en $\vect{v}$ in de figuur.\label{vb:h1-combi}}{%
\uitleg{$2\vect{u} - \vect{v} = 2\vect{u} + (-\vect{v})$. $\vect{u}$ gaat 2 rechts en 1 omhoog, dus $2\vect{u}$ gaat 4 rechts en 2 omhoog. $\vect{v}$ gaat 1 rechts en 2 omhoog, dus $-\vect{v}$ gaat 1 links en 2 omlaag.}
\naast{\begin{stappenvrij}{5cm}{2.6cm}
\text{teken } 2\vect{u} \text{ vanuit } P & 4 rechts, 2 omhoog. \\[1ex]
\text{teken } -\vect{v} \text{ aan de kop} & 1 links, 2 omlaag. \\[1ex]
\text{uitkomst: van } P \text{ naar de kop} & 3 rechts.
\end{stappenvrij}}{\begin{tikzpicture}[scale=.5,font=\scriptsize]
  \ruitjes{0}{0}{3}{3}
  \draw[vector] (0,1) -- (2,2) node[midway,above left,inner sep=1pt] {$\vect{u}$};
  \draw[vectorb] (2,0) -- (3,2) node[midway,right] {$\vect{v}$};
  \begin{scope}[xshift=4.5cm]
  \ruitjes{0}{0}{4}{3}
  \draw[vector] (0,1) -- (4,3) node[midway,above left,inner sep=1pt] {$2\vect{u}$};
  \draw[vectorb] (4,3) -- (3,1) node[midway,right] {$-\vect{v}$};
  \draw[som] (0,1) -- (3,1) node[midway,below] {$2\vect{u} - \vect{v}$};
  \fill (0,1) circle (2.5pt) node[left] {$P$};
  \end{scope}
\end{tikzpicture}}
De uitkomst is een horizontale vector van 3 hokjes naar rechts.}

Mag je bij vectoren ook haakjes uitwerken, zoals $2(\vect{u} + \vect{v}) = 2\vect{u} + 2\vect{v}$? In de figuur hieronder staat links de driehoek met $\vect{u}$, $\vect{v}$ en hun som. Rechts staan $2\vect{u}$, $2\vect{v}$ en $2\vect{u} + 2\vect{v}$. De grote driehoek is een vergroting met factor 2 van de kleine. Daarom is de derde zijde ook twee keer zo lang en evenwijdig: $2\vect{u} + 2\vect{v} = 2(\vect{u} + \vect{v})$.

\figuren{\begin{tikzpicture}[scale=.5,font=\scriptsize]
  \draw[vector] (0,0) -- (2,.5) node[midway,below right,inner sep=1pt] {$\vect{u}$};
  \draw[vectorb] (2,.5) -- (2.6,2) node[midway,right] {$\vect{v}$};
  \draw[som] (0,0) -- (2.6,2) node[midway,above left,inner sep=1pt] {$\vect{u} + \vect{v}$};
  \begin{scope}[xshift=5.5cm]
  \draw[vector] (0,0) -- (4,1) node[midway,below right,inner sep=1pt] {$2\vect{u}$};
  \draw[vectorb] (4,1) -- (5.2,4) node[midway,right] {$2\vect{v}$};
  \draw[som] (0,0) -- (5.2,4) node[midway,above left,inner sep=1pt] {$2\vect{u} + 2\vect{v}$};
  \end{scope}
\end{tikzpicture}}

\begin{theorie}
\begin{stelling}{Haakjes uitwerken}
Voor elk getal $k$ en alle vectoren $\vect{u}$ en $\vect{v}$ geldt

\vspace{1mm}
\qquad $\formule{k\,(\vect{u} + \vect{v}) = k\,\vect{u} + k\,\vect{v}}$
\end{stelling}
\end{theorie}

Met vectoren reken je dus zoals met letters in de algebra. Je mag de volgorde van een som wisselen, haakjes uitwerken, en gelijke vectoren samennemen: $3\vect{v} + 2\vect{v} = 5\vect{v}$, want vijf keer $\vect{v}$ na elkaar is vijf keer $\vect{v}$. Alleen een vector met een vector vermenigvuldigen doen we hier niet.

\voorbeeld{}{Vereenvoudig $3(\vect{u} - \vect{v}) + 2\vect{v} - \vect{u}$.\label{vb:h1-vereenvoudig}}{%
\uitleg{We werken de haakjes uit, en nemen dan de termen met $\vect{u}$ en de termen met $\vect{v}$ samen, zoals bij $3(a - b) + 2b - a$.}
\begin{stappen}
3(\vect{u} - \vect{v}) + 2\vect{v} - \vect{u} = 3\vect{u} - 3\vect{v} + 2\vect{v} - \vect{u} & Haakjes uitwerken. \\[1ex]
\phantom{3(\vect{u} - \vect{v}) + 2\vect{v} - \vect{u}} = 2\vect{u} - \vect{v} & $3\vect{u} - \vect{u}$ en $-3\vect{v} + 2\vect{v}$.
\end{stappen}
Dat is de vector die we in het vorige voorbeeld tekenden.}

\begin{oefeningen}
\oef{h1-combi}
\end{oefeningen}

\begin{alleenvoor}{4,5}
\needspace{16\baselineskip}
\subsection{Een eigenschap bewijzen met vectoren}

Met vectoren kun je ook eigenschappen van figuren bewijzen. Het idee: schrijf een vector op twee manieren, met Chasles en met veelvouden, en vergelijk. Een evenwijdigheid en een verhouding van lengtes volgen dan in één keer, want $\vect{w} = k\,\vect{v}$ zegt dat $\vect{w}$ evenwijdig is met $\vect{v}$ en $\abs{k}$ keer zo lang.

\voorbeeld{}{In driehoek $ABC$ is $M$ het midden van $[AB]$ en $N$ het midden van $[BC]$. Bewijs dat $MN$ evenwijdig is met $AC$, en half zo lang.\label{vb:h1-middenparallel}}{%
\uitleg{We schrijven $\vect{MN}$ met Chasles via $B$, en gebruiken dat $M$ en $N$ middens zijn. Dan zoeken we $\vect{AC}$ in de uitkomst.}
\naast{\begin{stappenvrij}{5.6cm}{3.8cm}
\vect{MN} = \vect{MB} + \vect{BN} & Chasles. \\[1ex]
\phantom{\vect{MN}} = \tfrac12\,\vect{AB} + \tfrac12\,\vect{BC} & $M$ en $N$ zijn middens. \\[1ex]
\phantom{\vect{MN}} = \tfrac12\,(\vect{AB} + \vect{BC}) & Haakjes. \\[1ex]
\phantom{\vect{MN}} = \tfrac12\,\vect{AC} & Chasles.
\end{stappenvrij}}{\begin{tikzpicture}[scale=.6,font=\scriptsize]
  \coordinate (A) at (0,0); \coordinate (B) at (1.6,3); \coordinate (C) at (5,0);
  \coordinate (M) at (.8,1.5); \coordinate (N) at (3.3,1.5);
  \fill[accent!8] (A) -- (B) -- (C) -- cycle;
  \draw[lijngrijs] (A) -- (B) -- (C);
  \draw[vector] (A) -- (C);
  \draw[som] (M) -- (N);
  \fill (M) circle (1.8pt) (N) circle (1.8pt);
  \node[below left] at (A) {$A$}; \node[above] at (B) {$B$}; \node[below right] at (C) {$C$};
  \node[left] at (M) {$M$}; \node[right] at (N) {$N$};
\end{tikzpicture}}
Omdat $\vect{MN} = \tfrac12\,\vect{AC}$, is $MN$ evenwijdig met $AC$, en $\lengte{MN} = \tfrac12\,\lengte{AC}$. Het lijnstuk $[MN]$ heet een \textit{middenparallel} van de driehoek.}

\end{alleenvoor}
%========================================================================================
\section{Krachten en snelheden}

Vectoren komen uit de fysica. Een kracht heeft een grootte (in newton, N), maar ook een richting en een zin: duwen of trekken, naar links of naar rechts. Ook een snelheid heeft een richting en een zin. Je tekent zo'n grootheid als een pijl \textit{op schaal}, bijvoorbeeld $1$~cm voor $10$~N. De lengte van de pijl toont de grootte.

\subsection{Krachten samenstellen}

Twee kinderen trekken met een touw aan een slee, elk in een andere richting. De slee beweegt alsof er één kracht aan trekt: de som van de twee krachten. Hoe groot is die kracht, en in welke richting trekt ze?

\begin{theorie}
Werken er twee of meer krachten op hetzelfde punt, dan is hun \textbf{\textcolor{accent}{resultante}} (of \textbf{\textcolor{accent}{resulterende kracht}}) de som van die krachten:

\vspace{1mm}
\qquad $\formule{\vect{F} = \vect{F}_1 + \vect{F}_2}$
\end{theorie}

De resultante heeft hetzelfde effect als de krachten samen. Omdat de krachten in hetzelfde punt aangrijpen, teken je de som meestal met de parallellogramregel.

\voorbeeld{}{Op een slee werken twee krachten die loodrecht op elkaar staan: $40$~N naar voor en $30$~N opzij. Teken de resultante op schaal en bepaal haar grootte.\label{vb:h1-kracht}}{%
\uitleg{Schaal: $1$~cm voor $10$~N, dus pijlen van 4~cm en 3~cm. Het parallellogram is hier een rechthoek. We meten de diagonaal, en controleren met de stelling van Pythagoras.}
\naast{\begin{stappenvrij}{5.6cm}{3.8cm}
\text{diagonaal: } 5 \text{ cm} & Gemeten. \\[1ex]
F = 5 \cdot 10 = 50 \text{ N} & Terug met de schaal. \\[1ex]
\sqrt{40^2 + 30^2} = \sqrt{2500} = 50 & Controle met Pythagoras.
\end{stappenvrij}}{\begin{tikzpicture}[scale=.55,font=\scriptsize]
  \fill[accent!8] (0,0) rectangle (4,3);
  \draw[hulplijn] (4,0) -- (4,3) -- (0,3);
  \draw[vector] (0,0) -- (4,0) node[midway,below] {$40$ N};
  \draw[vectorb] (0,0) -- (0,3) node[midway,left] {$30$ N};
  \draw[som] (0,0) -- (4,3) node[pos=.55,above left,inner sep=1pt] {$\vect{F}$};
  \rechtehoekxy{(4,0)}{(0,0)}{(0,3)}
\end{tikzpicture}}
De resultante is $50$~N. Ze is groter dan elk van de twee krachten, maar kleiner dan hun som $40 + 30 = 70$~N.}

Dat laatste is belangrijk: de grootte van de som is meestal niet de som van de groottes. Alleen als de krachten in dezelfde richting en zin werken, tel je de groottes gewoon op. Hoe groter de hoek tussen de krachten, hoe kleiner de resultante.

\voorbeeld{}{Twee krachten van elk $30$~N trekken aan een ring, met een hoek van $60^\circ$ tussen hen. Bepaal de resultante met een tekening op schaal.\label{vb:h1-kracht60}}{%
\uitleg{Schaal: $1$~cm voor $10$~N, dus twee pijlen van 3~cm met een hoek van $60^\circ$ (met de geodriehoek). We vullen aan tot een parallellogram (hier een ruit) en meten de diagonaal.}
\naast{\begin{stappenvrij}{5.4cm}{3.6cm}
\text{diagonaal: ongeveer } 5{,}2 \text{ cm} & Gemeten. \\[1ex]
F \approx 52 \text{ N} & Terug met de schaal.
\end{stappenvrij}}{\begin{tikzpicture}[scale=.75,font=\scriptsize]
  \coordinate (O) at (0,0); \coordinate (P) at (3,0); \coordinate (Q) at (60:3); \coordinate (R) at ($(P)+(Q)$);
  \fill[accent!8] (O) -- (P) -- (R) -- (Q) -- cycle;
  \draw[hulplijn] (P) -- (R) -- (Q);
  \draw[vector] (O) -- (P) node[midway,below] {$30$ N};
  \draw[vectorb] (O) -- (Q) node[midway,left] {$30$ N};
  \draw[som] (O) -- (R) node[pos=.7,below right,inner sep=1pt] {$\vect{F}$};
  \pic[draw,grijs,angle radius=4mm,pic text={$60^\circ$},angle eccentricity=1.9,pic text options={font=\tiny}] {angle=P--O--Q};
\end{tikzpicture}}
De resultante is ongeveer $52$~N, en niet $60$~N. Ze ligt precies midden tussen de twee krachten.}

\begin{werkwijze}
Kies een schaal. Teken de krachten vanuit één punt, met de juiste grootte en de juiste hoek. Teken de resultante met de parallellogramregel. Meet haar lengte en reken met de schaal terug naar newton. Staan de krachten loodrecht op elkaar, dan kun je met de stelling van Pythagoras ook rekenen.
\end{werkwijze}

Hangt een lamp stil aan twee kabels, dan is de resultante van alle krachten op de lamp de nulvector: de krachten heffen elkaar op. Dat heet \textit{evenwicht}. De resultante van de twee kabels is dan precies de tegengestelde van het gewicht.

\voorbeeld{}{Twee mensen trekken aan een kist met krachten van $40$~N naar rechts en $30$~N naar boven, zoals in voorbeeld~\ref{vb:h1-kracht}. Welke derde kracht moet een derde persoon uitoefenen om de kist stil te houden?\label{vb:h1-evenwicht}}{%
\uitleg{Bij evenwicht is $\vect{F}_1 + \vect{F}_2 + \vect{F}_3 = \nulvec$. Dus is $\vect{F}_3 = -(\vect{F}_1 + \vect{F}_2)$: de tegengestelde van de resultante.}
\naast{\begin{stappenvrij}{5.4cm}{3.6cm}
\vect{F}_1 + \vect{F}_2 = \vect{F}, \text{ met } F = 50 \text{ N} & Uit voorbeeld~\ref{vb:h1-kracht}. \\[1ex]
\vect{F}_3 = -\vect{F} & Even groot, andere zin.
\end{stappenvrij}}{\begin{tikzpicture}[scale=.45,font=\scriptsize]
  \draw[vector] (0,0) -- (4,0) node[midway,below] {$\vect{F}_1$};
  \draw[vectorb] (0,0) -- (0,3) node[midway,left] {$\vect{F}_2$};
  \draw[hulpvector] (0,0) -- (4,3) node[pos=.7,above left,inner sep=1pt] {$\vect{F}$};
  \draw[som] (0,0) -- (-4,-3) node[midway,below right,inner sep=1pt] {$\vect{F}_3$};
  \fill[inkt] (0,0) circle (3pt);
\end{tikzpicture}}
De derde persoon trekt met $50$~N, schuin naar links onder, in het verlengde van de resultante.}

\begin{oefeningen}
\oef{h1-kracht}
\end{oefeningen}

\subsection{Snelheden samenstellen}

Je wandelt op een rolpad in een luchthaven. Ten opzichte van de vloer ga je sneller dan wie naast het rolpad wandelt: de snelheid van het rolpad komt erbij. Ook snelheden tel je op als vectoren.

\voorbeeld{}{Een rolpad beweegt met $0{,}8$~m/s. Lina wandelt op het rolpad met $1{,}2$~m/s ten opzichte van het rolpad. Hoe snel beweegt ze ten opzichte van de vloer als ze met het rolpad mee wandelt? En als ze ertegenin wandelt?\label{vb:h1-rolpad}}{%
\uitleg{Alle snelheden liggen op één rechte. Wijzen ze naar dezelfde kant, dan liggen de pijlen in elkaars verlengde en tel je de groottes op. Wijzen ze naar elkaar toe, dan trek je af.}
\naast{\begin{stappenvrij}{4.6cm}{4cm}
1{,}2 + 0{,}8 = 2 \text{ m/s} & Mee: dezelfde zin. \\[1ex]
1{,}2 - 0{,}8 = 0{,}4 \text{ m/s} & Tegen: tegengestelde zin.
\end{stappenvrij}}{\begin{tikzpicture}[scale=.9,font=\scriptsize]
  \draw[vector] (0,.6) -- (1.2,.6); \draw[vectorb] (1.2,.6) -- (2,.6); \draw[som] (0,.3) -- (2,.3);
  \node[right] at (2.1,.45) {mee};
  \draw[vector] (0,-.25) -- (1.2,-.25); \draw[vectorb] (1.2,-.4) -- (.4,-.4); \draw[som] (0,-.65) -- (.4,-.65);
  \node[right] at (2.1,-.45) {tegen};
\end{tikzpicture}}
Wandelt ze tegen het rolpad in, dan gaat ze nog altijd vooruit, maar traag: met $0{,}4$~m/s in haar wandelrichting.}

Staan de snelheden niet op één rechte, dan tel je ze op met de kop-staartregel of de parallellogramregel.

\voorbeeld{}{Terug naar het probleem. Jens steekt met zijn kajak een rivier van 60~m breed over. Hij peddelt recht naar de overkant, met $1{,}5$~m/s ten opzichte van het water, maar het water stroomt met $0{,}8$~m/s en neemt hem mee. Waar komt hij aan de overkant aan, en hoe snel beweegt hij ten opzichte van de oever?\label{vb:h1-terug}}{%
\uitleg{De echte snelheid van Jens is de som van zijn eigen snelheid en die van het water. We tekenen de twee snelheden op schaal, kop aan staart. De som toont de richting waarin hij echt vaart: Jens volgt de gestreepte lijn in de figuur.}
\naast{\begin{stappenvrij}{6cm}{4.4cm}
\text{som: } 1{,}7 \text{ m/s} & Gemeten op de tekening. \\[1ex]
\sqrt{1{,}5^2 + 0{,}8^2} = \sqrt{2{,}89} = 1{,}7 & Controle met Pythagoras. \\[1ex]
60 : 1{,}5 = 40 \text{ s} & Elke seconde 1,5~m dichter bij de overkant. \\[1ex]
40 \cdot 0{,}8 = 32 \text{ m} & Zo ver drijft hij af.
\end{stappenvrij}}{\begin{tikzpicture}[scale=.7,font=\scriptsize]
  \fill[accent!8] (-.4,0) rectangle (2.4,3);
  \draw[grijs,thick] (-.4,0) -- (2.4,0) (-.4,3) -- (2.4,3);
  \draw[hulplijn] (0,0) -- (0,3) -- (1.6,3);
  \draw[accent,thick,dashed] (0,0) -- (1.6,3);
  \draw[vector] (0,0) -- (0,1.5) node[midway,left] {$1{,}5$};
  \draw[vectorb] (0,1.5) -- (.8,1.5) node[midway,above,inner sep=1pt] {$0{,}8$};
  \draw[som] (0,0) -- (.8,1.5);
  \node[right,accent] at (.5,.7) {$\vect{v}$};
  \fill[inkt] (0,0) circle (2pt) (1.6,3) circle (2pt);
  \draw[grijs,<->] (2.65,0) -- node[right] {60 m} (2.65,3);
  \draw[grijs,<->] (0,3.3) -- node[above] {32 m} (1.6,3.3);
\end{tikzpicture}}
Jens beweegt met $1{,}7$~m/s ten opzichte van de oever, schuin stroomafwaarts. Hij doet er 40 seconden over om de rivier over te steken, want de stroming verandert niets aan hoe snel hij de overkant nadert. In die tijd neemt het water hem 32~m mee: hij komt 32~m stroomafwaarts van het punt recht tegenover zijn vertrekpunt aan. Op een tekening van de rivier op schaal kun je dat punt ook vinden: trek een rechte vanuit het vertrekpunt in de richting van $\vect{v}$, tot aan de overkant.}

\begin{werkwijze}
Snelheden samenstellen: de snelheid ten opzichte van de grond is de som van de eigen snelheid (ten opzichte van het water of de lucht) en de snelheid van het water of de wind. Teken de pijlen op schaal, kop aan staart, en meet of reken de som.
\end{werkwijze}

\begin{oefeningen}
\oef{h1-snelheid}
\end{oefeningen}

%========================================================================================
\opdrachtenkop

\begin{oefeningen}
\oef{h1-o1}

\oef{h1-o2}

\oef{h1-o3}

\oef{h1-o4}

\oef{h1-o5}

\oef{h1-o6}

\oef{h1-o7}

\oef{h1-o8}

\oef{h1-o9}

\oef{h1-o10}

\oef{h1-o11}

\oef{h1-o12}

\oef{h1-o13}

\oef{h1-o14}

\oef{h1-o15}

\oef{h1-o16}

\oef{h1-o17}

\oef{h1-o18}

\oef{h1-o19}

\oef{h1-o20}

\oef{h1-o21}

\oef{h1-o22}

\oef{h1-o23}

\oef{h1-o24}
\end{oefeningen}

%----------------------------------------------------------------------------------------
\opdrachtsectie{Test jezelf}

\begin{oefeningen}
\oef{h1-t1}

\oef{h1-t2}

\oef{h1-t3}

\oef{h1-t4}

\oef{h1-t5}

\oef{h1-t6}

\oef{h1-t7}
\end{oefeningen}

\herhaalschema{%
  \fv{oef:h1-t1}{vb:h1-gelijk}{oef:h1-gelijk}%
  \fv{oef:h1-t2}{vb:h1-kopstaart}{oef:h1-kopstaart}%
  \fv{oef:h1-t3}{vb:h1-chasles}{oef:h1-chasles}%
  \fv{oef:h1-t4}{vb:h1-verschil2}{oef:h1-verschil}%
  \fv{oef:h1-t5}{vb:h1-combi}{oef:h1-combi}%
  \fv{oef:h1-t6}{vb:h1-vereenvoudig}{oef:h1-o6}%
  \fv{oef:h1-t7}{vb:h1-kracht}{oef:h1-kracht}}

\ifzelfstudie
\section*{Maak je eigen samenvatting}
Sluit dit boek en schrijf in je schrift: wat een vector is (richting, zin en grootte), wanneer twee vectoren gelijk zijn, de tegengestelde vector en de nulvector, de kop-staartregel en de parallellogramregel met een figuur, de betrekking van Chasles, hoe je een verschil tekent, wat $k\,\vect{v}$ betekent voor een positief en een negatief getal $k$, en hoe je de resultante van twee krachten bepaalt. Vergelijk daarna met de kaders.
\fi
